\chapter{Giant Cell Arteritis (Temporal Arteritis)}

\section*{Definition and Overview}
Giant cell arteritis (GCA), also called temporal arteritis, is a serious inflammatory disease of the large and medium arteries, particularly the branches of the carotid artery that supply the head, including the temporal artery. It causes the walls of affected arteries to become inflamed, narrowed, and sometimes blocked, reducing blood flow to the tissues they supply. GCA almost always affects people over the age of 50 and commonly causes new headaches, scalp tenderness, jaw pain when chewing, and visual disturbances. It is a medical emergency because it can cause permanent blindness if not treated urgently. Prompt treatment with high-dose corticosteroids usually prevents vision loss and relieves symptoms quickly.

\section*{Epidemiology}
Giant cell arteritis is the most common form of vasculitis in adults, affecting approximately 15--25 people per 100,000 over the age of 50 each year. It is rare before age 50 and becomes more common with each decade of life, peaking between ages 70 and 80. Women are affected about twice as often as men, and people of northern European descent are at higher risk. GCA frequently overlaps with polymyalgia rheumatica, a related inflammatory condition causing pain and stiffness of the shoulders and hips; around 40--50\% of people with GCA have polymyalgia rheumatica, and vice versa.

\section*{Aetiology and Pathophysiology}
The exact cause of giant cell arteritis is not fully understood, but it involves an abnormal immune response directed at the walls of arteries.

\begin{itemize}
    \item \textbf{Immune inflammation:} T-cells and macrophages invade the arterial wall, causing thickening, scarring, and the formation of the "giant cells" for which the disease is named
    \item \textbf{Vessel narrowing:} inflammation thickens the vessel wall and narrows the lumen, reducing blood flow
    \item \textbf{Thrombosis:} blood can clot within the narrowed vessel, further blocking flow
    \item \textbf{Ocular arteries:} blockage of the arteries to the optic nerve causes sudden, often permanent vision loss
    \item \textbf{Genetic and environmental triggers:} certain HLA genes increase susceptibility, and infections or other environmental factors may trigger the disease
    \item \textbf{Associated conditions:} GCA is closely related to polymyalgia rheumatica, and both respond to corticosteroids
\end{itemize}

\section*{Clinical Features}
The symptoms of GCA often develop gradually over weeks, but vision loss can occur suddenly and without warning.

\begin{itemize}
    \item New, persistent headache, often in the temple area, that is unlike previous headaches
    \item Scalp tenderness, such as pain when brushing the hair or resting the head on a pillow
    \item Jaw claudication: pain or fatigue in the jaw when chewing
    \item Visual symptoms: blurred vision, double vision, or transient vision loss, which can progress to permanent blindness
    \item Fever, fatigue, weight loss, and general malaise
    \item A tender, thickened, or pulseless temporal artery
    \item Pain and stiffness of the shoulders and hips when polymyalgia rheumatica coexists
    \item Rarely, symptoms of large-vessel involvement such as arm or leg claudication
\end{itemize}

\section*{Differential Diagnosis}
Several conditions can mimic the symptoms of GCA and must be considered.

\begin{itemize}
    \item Tension headache, migraine, and other primary headache disorders
    \item Temporal artery pathology such as aneurysm or dissection
    \item Other vasculitides, including polyarteritis nodosa and ANCA-associated vasculitis
    \item Infection, including shingles affecting the scalp or face
    \item Dental and temporomandibular joint problems causing jaw pain
    \item Cervical spine disease causing neck and occipital pain
    \item Cancers presenting with fever, weight loss, and headache
    \item Fibromuscular dysplasia and atherosclerosis of the carotid arteries
\end{itemize}

\section*{When to Seek Medical Care}
Anyone over 50 with a new headache, scalp tenderness, jaw pain on chewing, or visual symptoms should see a doctor urgently, as early treatment prevents blindness. Do not wait for symptoms to resolve on their own.

\textbf{Call 000 (or local emergency services) for:}
\begin{itemize}
    \item Sudden vision loss or profound visual disturbance
    \item Severe sudden headache, especially with neurological symptoms, which may indicate other emergencies such as stroke or subarachnoid haemorrhage
    \item Chest pain, breathlessness, or signs of a heart attack or stroke, which can occur with large-vessel involvement
\end{itemize}

\section*{Diagnosis and Investigations}
Because blindness can occur rapidly, treatment may begin before tests confirm the diagnosis.

\begin{itemize}
    \item \textbf{Clinical assessment:} the combination of age, new headache, scalp tenderness, jaw claudication, and raised inflammatory markers strongly suggests GCA
    \item \textbf{Blood tests:} erythrocyte sedimentation rate (ESR) and C-reactive protein (CRP) are markedly elevated in most cases
    \item \textbf{Full blood count:} may show anaemia and a high platelet count
    \item \textbf{Temporal artery biopsy:} a small sample of the temporal artery is examined under a microscope for giant cells and inflammation; this is the gold standard, though treatment is started first
    \item \textbf{Ultrasound:} high-resolution ultrasound of the temporal and axillary arteries can show the "halo sign" of vessel wall inflammation
    \item \textbf{Imaging of large vessels:} MRI, CT, or PET can identify involvement of the aorta and its branches
\end{itemize}

\section*{Management}
Treatment must begin immediately when GCA is suspected, without waiting for biopsy or imaging results.

\begin{itemize}
    \item \textbf{High-dose corticosteroids:} prednisone or prednisolone is started at once; symptoms usually improve within days and vision loss is prevented
    \item \textbf{Long-term taper:} the dose is gradually reduced over many months, guided by symptoms and inflammatory markers
    \item \textbf{Bone protection:} calcium, vitamin D, and sometimes bisphosphonates are used to protect against steroid-induced osteoporosis
    \item \textbf{Monitoring for relapse:} regular review is needed because symptoms can recur as the dose is lowered
    \item \textbf{Steroid-sparing medicines:} tocilizumab, an immune-modifying medicine, may be used to reduce the steroid dose needed
    \item \textbf{Low-dose aspirin:} sometimes recommended to reduce the risk of clotting complications
    \item \textbf{Referral:} all patients are managed by a rheumatologist, often in collaboration with an ophthalmologist
\end{itemize}

\section*{Complications}
\begin{itemize}
    \item Permanent blindness, which occurs in up to 15--20\% of untreated or undertreated patients
    \item Stroke, particularly involving the posterior circulation
    \item Aortic aneurysm and dissection, which can occur years later with large-vessel disease
    \item Complications of corticosteroid therapy, including osteoporosis, diabetes, weight gain, infection, and hypertension
    \item Relapse, which occurs in around half of patients during steroid tapering
    \item Polymyalgia rheumatica symptoms in those with overlap disease
\end{itemize}

\section*{Prognosis and Outcomes}
With prompt and adequate treatment, the outlook for GCA is good. Symptoms settle quickly, and the risk of blindness falls dramatically once corticosteroids are started. Most patients require one to two years of treatment, and many can eventually stop all medication, though relapse is common during the taper. The main long-term concerns are the side effects of prolonged steroid use and the small but real risk of late aortic complications, so ongoing specialist follow-up is recommended. Untreated, GCA can cause permanent blindness and carries a real risk of stroke.

\section*{Prevention and Self-Care}
\begin{itemize}
    \item Report any new headache, scalp tenderness, jaw pain on chewing, or visual changes promptly, especially if you are over 50
    \item Do not delay medical review while waiting for symptoms to settle
    \item Take corticosteroids exactly as prescribed and never stop them suddenly
    \item Attend all follow-up appointments so the dose can be tapered safely
    \item Protect your bones with calcium and vitamin D, and discuss bone density screening with your doctor
    \item Monitor and report any new symptoms such as chest pain, or arm and leg claudication, which may indicate large-vessel disease
    \item Maintain a healthy lifestyle, including a balanced diet and regular activity, to counter steroid side effects
\end{itemize}

\section*{Sources and Further Reading}
The following reputable sources provide further information for healthcare students and patients seeking reliable, current advice about giant cell arteritis.

\begin{itemize}
    \item Australian Rheumatology Association. \textit{Giant cell arteritis information for patients.} https://rheumatology.org.au/
    \item Healthdirect Australia. \textit{Giant cell arteritis (temporal arteritis).} https://www.healthdirect.gov.au/
    \item Vasculitis Foundation. \textit{Giant cell arteritis: symptoms, diagnosis, and treatment.} https://www.vasculitisfoundation.org/
    \item NHS. \textit{Giant cell arteritis.} https://www.nhs.uk/
    \item Mayo Clinic. \textit{Giant cell arteritis: symptoms and causes.} https://www.mayoclinic.org/
    \item MSD Manual. \textit{Giant cell arteritis.} https://www.msdmanuals.com/
\end{itemize}
